Can restoring evidence reverse unsupported approval learned through peer
interaction? We tested this recovery question in a synthetic review task with
exact answers. Across five independent paired blocks, a model reviewer interacted
with scripted peers that always accepted or always rejected. Tasks, restricted
evidence, and neutral messages were matched; approving correct claims earned
points even without the evidence required by the verification rule. We then
tested evidence restoration, explicit correction, and correction combined with
history removal on new correct and faulty claims. The completed study produced
340 scored decisions. Neither peer condition produced unsupported approval:
there were 0/45 approvals per condition in the prespecified formation comparison.
Each of the five full-evidence conditions rejected all 20 faulty claims and
approved all 20 correct claims. The restricted-evidence condition rejected all
40 claims. The primary false-approval difference between peer histories was
zero, as were the observed added differences from correction and history removal.
Identical outcomes across five blocks did not yield informative confidence
intervals. Because no peer-induced failure was demonstrated, these results do
not establish recovery from collusion or the effectiveness of correction and
reset. They document protocol-consistent verdicts in this configuration and leave
the intended recovery question open under conditions that establish the failure
to be repaired.
